\documentclass[11pt]{article}
% Shared local declarations. Document roots use the repository's business-admin family.
\usepackage{business-admin}
\usepackage{amsmath,amssymb,tabularx,longtable,array,booktabs,enumitem}
\usepackage{ragged2e}
\newcolumntype{Y}{>{\RaggedRight\arraybackslash}X}
\newcommand{\SuiteMeta}[2]{\def\DocTitle{#1}\def\DocSubtitle{#2}\def\DocVersion{1.0}\def\DocStatus{Educational and operational reference}\def\DocOwner{Parliamentary Procedure and Governance Suite}\def\DocSummary{Independently authored reference and working materials, based on RONR 12th edition.}}
\newcommand{\SuiteOpen}{\makedoctitle\vspace{-0.5em}\noindent\rule{\linewidth}{0.4pt}\par\medskip}
\newcommand{\RonrAnchor}[1]{\par\smallskip{\footnotesize\textbf{Primary reading:} Robert's Rules of Order Newly Revised, 12th ed., #1. Original synthesis; not official interpretation.}\par}
\newcommand{\WriteLine}{\par\noindent\rule{\linewidth}{0.35pt}\par\vspace{0.5cm}}
\newcommand{\WriteLines}[1]{\par\noindent\rule{\linewidth}{0.35pt}\par\vspace{0.35cm}\noindent\rule{\linewidth}{0.35pt}\par\vspace{0.35cm}\noindent\rule{\linewidth}{0.35pt}\par\vspace{0.25cm}}
\newcommand{\Field}[1]{\par\medskip\noindent\textbf{#1}\par\WriteLines{3}}
\newcommand{\GovCheckbox}{\(\square\)}
\setlength{\emergencystretch}{3em}

\SuiteMeta{Bylaws, Formation and Change Control}{Technical reference | original application synthesis}
\begin{document}
\SuiteOpen
\section{Scope and interpretation}
This is an original operational reference based on RONR (12th ed.). It is subordinate to applicable law and the organization’s validly adopted governing documents. Rule citations identify the primary source; the authors of this suite have not issued official interpretations.

\section{Foundational design}
Bylaws should state organizational purpose, membership, officers, elections, meetings, quorum, boards, committees where required, parliamentary authority and amendment process. Separate durable governance provisions from easily revised routine arrangements.

\textbf{Practical audit hook.} Source anchor: RONR 53--57.

\section{Formation sequence}
Mass meetings and permanent societies have different organizational tasks. A new permanent society generally needs a structured organizing process before exercising permanent powers.

\textbf{Practical audit hook.} Lab: produce first-meeting minutes, bylaws committee referral and second-meeting adoption records.

\section{Amendment review}
Test previous notice, amendment scope, threshold, voting body, required external approvals and effective date; do not substitute fictional defaults for the organization’s own text. The 12th edition includes a specific limitation on who may prepare a bylaws revision (57:5).

\textbf{Practical audit hook.} Assurance: report overbroad proposed amendments and revised language.

\section{Legal boundary}
Mergers, dissolution, property transfers, membership and officer duties may be governed by jurisdiction-specific law and filings beyond parliamentary procedure. Engage qualified legal counsel before operating changes.

\textbf{Practical audit hook.} Document: external-law review with named reviewer and source date.

\section{Independent review checklist}
\begin{enumerate}[leftmargin=*,itemsep=4pt]
\item Find the controlling adopted text and note its revision date.
\item Specify the decision, eligible decision maker, notice, quorum and vote basis.
\item Record evidence without fabricating approvals or preserving private member data in a public repository.
\item Classify unresolved issues as parliamentary, legal, governance policy or factual disputes.
\end{enumerate}
\end{document}
